% Created 2026-09-08 Tue 08:27
% Intended LaTeX compiler: pdflatex
\documentclass{beamer}\usepackage{listings}
\usepackage{color}
\usepackage{amsmath}
\usepackage{array}
\usepackage[T1]{fontenc}
\usepackage{natbib}
\lstset{
keywordstyle=\color{blue},
commentstyle=\color{red},stringstyle=\color[rgb]{0,.5,0},
basicstyle=\ttfamily\small,
columns=fullflexible,
breaklines=true,
breakatwhitespace=false,
numbers=left,
numberstyle=\ttfamily\tiny\color{gray},
stepnumber=1,
numbersep=10pt,
backgroundcolor=\color{white},
tabsize=4,
literate={~}{$\sim$}{1}{ø}{{\o}}1{æ}{{\ae}}1{å}{{\aa}}1{Ø}{{\OE}}1{Æ}{{\AE}}1{Å}{{\AA}}1,keepspaces=true,
showspaces=false,
showstringspaces=false,
xleftmargin=.23in,
frame=single,
basewidth={0.5em,0.4em},
}
\institute{}
\subtitle{}
\RequirePackage[absolute, overlay]{textpos}
\setbeamertemplate{footline}[frame number]
\setbeamertemplate{navigation symbols}{}
\RequirePackage{fancyvrb}
\RequirePackage{array}
\RequirePackage{multirow}
\RequirePackage{tcolorbox}
\definecolor{mygray}{rgb}{.95, 0.95, 0.95}
\newcommand{\mybox}[1]{\vspace{.5em}\begin{tcolorbox}[boxrule=0pt,colback=mygray] #1 \end{tcolorbox}}
\newcommand{\sfootnote}[1]{\renewcommand{\thefootnote}{\fnsymbol{footnote}}\footnote{#1}\setcounter{footnote}{0}\renewcommand{\thefootnote}{\arabic{foot note}}}
\usepackage{alphalph}
\alphalph{\value{footnote}}
\setbeamertemplate{itemize item}{\textbullet}
\setbeamertemplate{itemize subitem}{-}
\setbeamertemplate{itemize subsubitem}{}
\setbeamertemplate{enumerate item}{\insertenumlabel.}
\setbeamertemplate{enumerate subitem}{\insertenumlabel.\insertsubenumlabel}
\setbeamertemplate{enumerate subsubitem}{\insertenumlabel.\insertsubenumlabel.\insertsubsubenumlabel}
\setbeamertemplate{enumerate mini template}{\insertenumlabel}
\makeatletter\def\blfootnote{\xdef\@thefnmark{}\@footnotetext}\makeatother
\newcommand{\E}{\ensuremath{\mathrm{E}}}
\renewcommand{\P}{\ensuremath{\mathrm{P}}}
\renewcommand{\d}{\ensuremath{\mathrm{d}}}
\setbeamertemplate{caption}{\raggedright\insertcaption\par}
\usepackage{amsmath}
\DeclareMathOperator*{\argmin}{arg\,min}
\DeclareMathOperator*{\argmax}{arg\,max}

\renewcommand*\familydefault{\sfdefault}
\itemsep2pt
\usepackage[utf8]{inputenc}
\usepackage[T1]{fontenc}
\usepackage{graphicx}
\usepackage{longtable}
\usepackage{wrapfig}
\usepackage{rotating}
\usepackage[normalem]{ulem}
\usepackage{amsmath}
\usepackage{amssymb}
\usepackage{capt-of}
\usepackage{hyperref}
\usetheme{default}
\author{Hjalte Søberg Mikkelsen}
\date{}
\title{}
\begin{document}

\section{Introduction}
\label{sec:orga5f6a42}
\begin{frame}[label={sec:org581a002}]{Introduction}
My PhD project is part of a larger project called DP-Next. DP-Next is a project aimed at developing a sustainably effective strategy for prevention of Type 2 Diabetes in Denmark, Greenland and the Faroe Islands.
\vfill
My goal is to develope a risk prediction model that can predict the 5 year risk of getting diagnosed with type 2 diabetes and the 5 year risk of death, only using registry data. This model is to be used to select people that will be included in the devolopement of the prevention strategy. 
\end{frame}


\begin{frame}[label={sec:orgeed6438}]{Motivation/Setting}
When you estimate the cause specific risks at time \(\tau\) using semi-parametric Cox regression for the cause specific hazards, the sum of all cause specific risks can be more than 1 for certain covariate values.
\vfill
We need to export the prediction model developed in protected environments, such that we can give it to the public.
\vfill

It is unclear when to artifical censor subjects at the prediction time horizon.
\end{frame}

\begin{frame}[label={sec:org89b2e31}]{Setting: competing risks}
Assume we are in a competing risks scenario and want to predict the risk of getting diagnosed with T2D and the risk of dying, at some fixed timepoint \(\tau\). Denote:
\vfill
\begin{itemize}
\item \(T\) as the time between time0 and the date of an event
\item \(D\in\{1,2\}\) as the cause of the event
\item \(X = (X^1,...,X^p)\) as a p-dimensional vector of baseline covariates, which can both be numeric or categorical.
\end{itemize}
\vfill
Additionally we assume that the event time \(T\) is right censored at the random time \(C\), where \(P(C>\tau\vert X)>0\) almost surely. Denote \(\tilde{T} = \min(T,C)\) and \(\tilde{D} = \Delta D\) as the observed time and event, where \(\Delta = 1\{T\leq C\}\).
\end{frame}

\begin{frame}[label={sec:orgbfbee36}]{Setting: Proportional hazards model}
Let \(\beta\) be a p-dimensional parameter vector and \(x\) be some observed value of \(X\). We are interrested in the cumulative incidence function for both causes, which for cause \(j\) is \(F_j(t|x) = P(T\leq t, D = j|x)\), and can be written as
\begin{align*}
F_j(t|x) = \int_0^t S(s-|x)d\Lambda_j(s|x),
\end{align*}
where \(S\) is the survival function
\begin{align*}
S(t|x) = \exp\left(-\sum_{i = 1}^2 \Lambda_i(s|x)\right),
\end{align*}
and \(\Lambda_j\) is the cumulative hazard
\begin{align*}
\Lambda_j(t|x) = \exp(x\beta_j)\Lambda_{0j}(t).
\end{align*}
\(\Lambda_{0j}\) is the baseline cumulative hazard for cause j, which is defined as \(\Lambda_{0j}(t) = \int_0^t \lambda_{0j}(s)ds\), where \(\lambda_{0j}\) is the baseline hazard
\end{frame}

\begin{frame}[label={sec:org9f663b4}]{Proportional hazards model estimation}
The typical ways of etimating these kinds of models is to either use partial likelihood, where we assume the baseline hazard is non-parametric, or maximum likelihood, where we need to choose some parametric form for the baseline hazard functions. \pause
\vfill
In this project I will choose a parametric form for the baseline hazards, but I will minimize the Brier-score at a fixed timepoint instead of maximizing the likelihood.\pause
\vfill
I will throughout the presentation assume that \(C\) and \(T\) are conditionally independent given \(X\).
\end{frame}

\begin{frame}[label={sec:orgfc1ea4b}]{Maximum likelihood}
Assume we have \(N\) observations of \((\tilde{T},\tilde{D}, X)\). Using that the survival function is a product of the individual cause-specific survival functions, \(S(t_i|x_i) = \prod S_j(t_i|x_i)\), we can write the likelihood as
\begin{align*}
L = \prod_{j = 1}^2\prod_{i = 1}^N\lambda_{j}(t_i|x_i)^{d_{ij}}\exp(-\Lambda_j(t_i|x_i)),
\end{align*}
where \(d_{ij} = 1\{\tilde{D}_i = j\}\). \pause This means that the likelihood function is a product of 2 likelihoods, one for each cause. We can then maximize each cause-specific likelihood function seperately by choosing some parametric form for the baseline hazards.
\end{frame}


\begin{frame}[label={sec:orgb978128}]{Brier score estimation}
The model I will be focusing on is the Cox-exponential model. This model assumes that the baseline hazard is constant, meaning that
\begin{align*}
\Lambda_{0j}(\tau) = \int_0^{\tau} \lambda_{0j}(s)ds = \tau\lambda_{0j},
\end{align*}
as \(\lambda_{0j}\) is now a constant. \pause

I will then estimate the parameters by minimizing the Brier-score for some fixed timepoint \(\tau\), i.e.

\begin{align*}
BSE = &\argmin_{\boldsymbol{\beta}\in\mathbb{R}^{3\times p},\boldsymbol{\lambda_0}\in\mathbb{R}^{+3}}\sum_{j = 0}^2\frac{1}{N}\sum_{i = 1}^N(F_j(\tau|x_i,\boldsymbol{\beta},\boldsymbol{\lambda_0})-o_i(\tau,j))^2,
\end{align*}
where \(o_i(\tau,j) = 1(\tilde{D}_i= j,\tilde{T}_i\leq \tau)\). Note that we also estimate the censoring state (Munch and Gerds, 2024). \pause

Assuming constant baseline hazards for all causes and letting \(A_i = \sum_{k = 0}^2 \exp(\boldsymbol{\beta}_kx_i)\lambda_{0k}\), we can write the BSE as
\begin{align*}
BSE = \argmin_{\boldsymbol{\beta}\in\mathbb{R}^{3\times p},\boldsymbol{\lambda_0}\in\mathbb{R}^{+3}}\sum_{j = 0}^2\sum_{i = 1}^N \left(\exp(\boldsymbol{\beta}_jx_i)\lambda_{0j}\frac{1-\exp\left[-\tau A_i\right]}{A_i}-o_i(\tau,j)\right)^2
\end{align*}
\end{frame}




\begin{frame}[label={sec:org4cf67c7}]{Artificial censoring}
Asume we will artificial censor at \(t = 5\). This means our new observed time will now be \(\tilde{T}_{\text{trunc}} = \min(\tilde{T},5) = \min(T,C,5)\) and the new observed event \(\tilde{D}_{\text{trunc}} = 1\{T\leq 5\}\tilde{D} =  1\{T\leq C \wedge T\leq 5\}D\). Doing artificial censoring means we truncate all times at a specific timepoint, and practically means that estimations methods such as the partial likelihood and maximum likelihood can't learn from data after that timepoint. The Brier-score estimation is unaffected by artificial censoring, as it only learns from data before timepoint \(\tau\).
\end{frame}

\begin{frame}[label={sec:org75a1911}]{Artificial censoring example plot using non-proportional hazards data}
\end{frame}




\begin{frame}[label={sec:orgfe5f16e}]{Practical examples}
Using DP-Next data estimate for specific covariate values and fit models both with and without artificial censoring
\end{frame}


\begin{frame}[label={sec:org1273f6c}]{Discussion}
The BSE is not convex so it needs good initial guess to perform well.
\vfill
\end{frame}
\end{document}
